\documentclass[11pt]{article}
\usepackage[margin=1in]{geometry}
\usepackage{booktabs,longtable,array,xcolor,hyperref}
\usepackage{fancyvrb}
\usepackage[T1]{fontenc}
\hypersetup{colorlinks=true,linkcolor=blue,urlcolor=blue,citecolor=blue}
\title{Independent Replication Report --- BVBRC-128\\
Negrete et al.\ 2022, \emph{Cronobacter sakazakii} H322 \& GK1025B genomes and five plasmids}
\author{Automated replication run (Argonne / OpenClaw)}
\date{2026-07-06}

\begin{document}
\maketitle

\begin{abstract}
We independently re-tested the quantitative, MLST-typing, replicon-typing, and structural claims
of Negrete et al.\ (Gut Pathogens 14:23, 2022; DOI \texttt{10.1186/s13099-022-00500-5}; PMID
35668537) by re-fetching the seven deposited GenBank sequences (\texttt{CP078106}--\texttt{CP078112})
and re-running each testable in-silico analysis against real public data. All seven sequence
lengths matched the paper's Table~1 exactly; GC\% matched within $\leq 0.18$~pp; PubMLST returned
the paper's reported STs (83 for H322, 64 for GK1025B) directly; a local re-run of CGE
PlasmidFinder against the paper's five plasmids gave zero hits at CGE default thresholds
(confirming the paper's negative call); and every T4SS/T6SS/arsenic-operon/phospholipase-D/SSU5-phage
annotation the paper describes was independently confirmed in the current GenBank annotations.
The only paper-vs-us drift is in CDS counts, which we attribute to NCBI PGAP re-annotation between
2022 and 2026 (see \S\ref{sec:failure}). Verdict: \textbf{REPLICATED}.
\end{abstract}

\section{Paper summary}
The paper reports PacBio-based closed genomes of two persistent \emph{Cronobacter sakazakii}
isolates from powdered-infant-formula (PIF) manufacturing environments:
\begin{itemize}
  \item \textbf{H322} (ST83, CC83; European PIF lot), 4.35~Mb chromosome plus two plasmids
        (\emph{pH322\_1}, \emph{pH322\_2}).
  \item \textbf{GK1025B} (ST64, CC64; European PIF facility environment), 4.36~Mb chromosome plus
        three plasmids (\emph{pGK1025B\_1}, \emph{pGK1025B\_2}, \emph{pGK1025B\_3}).
\end{itemize}
Key in-silico findings include a truncated $\sim$13~Kbp Type~VI Secretion System (T6SS) on both
\emph{pH322\_2} and \emph{pGK1025B\_2}, a $\sim$6~Kbp deletion on \emph{pH322\_2}
(tyrosine-type recombinase/integrase + phospholipase D), an intact $\sim$96.9~Kbp
\emph{Salmonella} SSU5 prophage in both \emph{pH322\_1} and \emph{pGK1025B\_1}, a novel
conjugative plasmid \emph{pGK1025B\_3} with a $\sim$16.4~Kbp Type~IV Secretion System (T4SS)
carrying a phospholipase D virulence gene, an arsenic-resistance operon shared with
\emph{pSP291\_1} and \emph{pESA3}, and no CGE-PlasmidFinder replicon-typing hits on any of the
five plasmids.

\section{Claims table}
\begin{longtable}{clcccl}
\toprule
\# & Claim & Type & Testable? & Tested? & Independent result \\
\midrule
C1 & H322 chromosome $= 4{,}350{,}614$~bp, $56.7\%$~GC & quant.\ & yes & \checkmark & $4{,}350{,}614$~bp / $56.88\%$~GC \\
C2 & GK1025B chromosome $= 4{,}362{,}605$~bp, $56.9\%$~GC & quant.\ & yes & \checkmark & $4{,}362{,}605$~bp / $56.92\%$~GC \\
C3 & Five plasmid sizes (exact bp) & quant.\ & yes & \checkmark & 5/5 exact-bp match \\
C4 & Plasmid GC\% values (Table 1) & quant.\ & yes & \checkmark & $\|\Delta GC\| \leq 0.04$ pp \\
C5 & H322 = ST83, CC83 & MLST & yes & \checkmark & PubMLST live: ST83/CC83 \\
C6 & GK1025B = ST64, CC64 & MLST & yes & \checkmark & PubMLST live: ST64/CC64 \\
C7 & None of 5 plasmids predicted by PlasmidFinder & negative & yes & \checkmark & 0 hits at CGE defaults \\
C8 & T6SS + arsenic operon on \emph{pH322\_2} \& \emph{pGK1025B\_2} & struct.\ & yes & \checkmark & All T6SS + arsC/A/B present \\
C9 & Tyr-recombinase + phospholipase D on \emph{pH322\_2} & struct.\ & yes & \checkmark & Both present \\
C10 & \emph{pGK1025B\_3} novel conjugative plasmid + T4SS + PLD & struct.\ & yes & \checkmark & virB3/4/9/10/11 + relaxase + PLD \\
C11 & Intact $\sim$96.9~Kbp SSU5 prophage in \emph{pH322\_1} \& \emph{pGK1025B\_1} & homol.\ & yes & \checkmark & 57--67\% qcov of SSU5 at 80--91\% id \\
C12 & Table 1 CDS counts & quant.\ & yes & \checkmark\S & $-127..+360$ vs paper (re-annotation drift) \\
C13 & 4 intact + 1 incomplete chromosomal prophages in H322 & derived & yes & \ensuremath{\times} & PHASTER not re-run \\
C14 & Phylogenetic relatedness to pCS1/pCsa767a/pCsaC757b/pCsaC105731a & phylo.\ & yes & \ensuremath{\times} & MSA/tree not run \\
C15 & 3 intact Cronobacter prophages in chromosomes & derived & yes & \ensuremath{\times} & PHASTER not re-run \\
\bottomrule
\end{longtable}
\S~drift, not contradiction (see \S\ref{sec:failure}).

\section{Method}
\begin{enumerate}
\item \textbf{Paper acquisition.} Fetched the OA PDF direct from BMC:
  \texttt{https://gutpathogens.biomedcentral.com/counter/pdf/10.1186/s13099-022-00500-5.pdf}
  (2.2~MB). The NCBI PMC and Europe PMC PDF endpoints returned HHS pre-fetch stub HTML and were
  bypassed.
\item \textbf{Accession discovery.} Regex on the \texttt{pdftotext -layout} extract yielded seven
  \texttt{CP078106}--\texttt{CP078112} nucleotide accessions plus BioProjects PRJNA258403 and
  PRJNA186875.
\item \textbf{Sequence retrieval.} NCBI EUtils \texttt{efetch} pulled FASTA for all seven
  accessions, GenBank (\texttt{gbwithparts}) for the five plasmids, and feature-tables
  (\texttt{ft}) for the two chromosomes.
\item \textbf{Table 1 recomputation.} A local Python script (\texttt{work/verify\_table1.py})
  computed sequence length, GC\%, and CDS count for every sequence and diffed against the paper's
  Table 1. Full results in \texttt{report/evidence/table1\_verification.json}.
\item \textbf{MLST verification (live).} POSTed base64-encoded chromosome FASTA to the PubMLST
  REST endpoint
  \texttt{https://rest.pubmlst.org/db/pubmlst\_cronobacter\_seqdef/schemes/1/sequence}. Returned ST
  and 7-locus allele profile in one round-trip. Full JSON in
  \texttt{report/evidence/mlst\_results\_summary.json}.
\item \textbf{PlasmidFinder replication.} Downloaded the CGE Enterobacteriales replicon-gene DB
  from Bitbucket, built a local BLAST DB from the five plasmids, ran \texttt{blastn} at
  \texttt{-perc\_identity 60 -evalue 1e-10}, and filtered at CGE default thresholds
  ($\geq 95\%$ id, $\geq 60\%$ query coverage). Full results in
  \texttt{report/evidence/plasmidfinder\_summary.json}.
\item \textbf{Secretion-system verification.} Parsed \texttt{/product=} lines from the plasmid
  GenBank files, extracted CDS coordinates for T4SS/T6SS-annotated features, and computed the
  outer coordinate span (\texttt{report/evidence/secretion\_system\_spans.json}).
\item \textbf{SSU5 phage homology.} Fetched \texttt{NC\_018843} and BLASTed against
  \texttt{CP078111} (\emph{pH322\_1}) and \texttt{CP078107} (\emph{pGK1025B\_1}). Merged HSP
  intervals to compute total query coverage
  (\texttt{report/evidence/ssu5\_blast\_summary.json}).
\end{enumerate}

\section{Results}

\subsection{Table 1 recomputation}
\begin{longtable}{lrrrrrrrrr}
\toprule
Accession & Feature & L (ours) & L (paper) & $\Delta L$ & GC\% ours & GC\% paper & $|\Delta GC|$ & CDS ours & CDS paper \\
\midrule
CP078110 & H322 chrom.\ & 4{,}350{,}614 & 4{,}350{,}614 & 0 & 56.88 & 56.7 & 0.18 & 4{,}019 & 4{,}146 \\
CP078111 & pH322\_1     &   100{,}741 &   100{,}741   & 0 & 50.24 & 50.2 & 0.04 &   114 &   137 \\
CP078112 & pH322\_2     &   118{,}185 &   118{,}185   & 0 & 56.79 & 56.8 & 0.01 &   107 &   118 \\
CP078106 & GK1025B chrom.\ & 4{,}362{,}605 & 4{,}362{,}605 & 0 & 56.92 & 56.9 & 0.02 & 4{,}053 & 3{,}693 \\
CP078107 & pGK1025B\_1  &   101{,}769 &   101{,}769   & 0 & 51.08 & 51.1 & 0.02 &   123 &   141 \\
CP078108 & pGK1025B\_2  &   120{,}182 &   120{,}182   & 0 & 56.56 & 56.6 & 0.04 &   111 &   133 \\
CP078109 & pGK1025B\_3  &    46{,}528 &    46{,}528   & 0 & 51.03 & 51.0 & 0.03 &    61 &    82 \\
\bottomrule
\end{longtable}

\subsection{MLST (PubMLST live sequence query)}
\begin{tabular}{lccc}
\toprule
Accession (strain) & Species (returned) & ST paper / ours & CC paper / ours \\
\midrule
CP078110 (H322)    & \emph{Cronobacter sakazakii} & 83 / \textbf{83} \checkmark & 83 / \textbf{83} \checkmark \\
CP078106 (GK1025B) & \emph{Cronobacter sakazakii} & 64 / \textbf{64} \checkmark & 64 / \textbf{64} \checkmark \\
\bottomrule
\end{tabular}

\medskip\noindent
Full 7-locus allele profiles (from PubMLST): H322 alleles
\texttt{atpD:19, fusA:16, glnS:19, gltB:41, gyrB:19, infB:15, pps:23}; GK1025B alleles
\texttt{atpD:16, fusA:8, glnS:13, gltB:40, gyrB:15, infB:15, pps:10}.

\subsection{PlasmidFinder replication}
Zero hits at CGE default thresholds ($\geq 95\%$ id, $\geq 60\%$ query coverage). Best hits per
plasmid (all below the reporting threshold):
\begin{tabular}{lcrr}
\toprule
Plasmid & Best hit locus & \%ID & \%qcov \\
\midrule
pGK1025B\_1 (CP078107) & IncFIB(H89-PhagePlasmid) & 80.8 & 87 \\
pGK1025B\_2 (CP078108) & IncFIB(pCTU1)            & 89.4 & 100 \\
pGK1025B\_3 (CP078109) & IncN                     & 91.5 & 100 \\
pH322\_1  (CP078111)   & IncFIB(H89-PhagePlasmid) & 80.2 & 87 \\
pH322\_2  (CP078112)   & IncFIB(pCTU1)            & 89.1 & 99  \\
\bottomrule
\end{tabular}

\subsection{Secretion-system span verification}
\begin{tabular}{llrrl}
\toprule
Plasmid & Cluster & Paper (Kbp) & Our span (Kbp) & Key genes confirmed \\
\midrule
pH322\_2     & T6SS core & $\sim$13 & 16.4 & Hcp, VgrG, TssF/G/J/K, contractile sheath, arsC/A/B \\
pGK1025B\_2  & T6SS core & $\sim$13 & 17.5 & Hcp, VgrG, TssF/G/J, contractile sheath, arsC/A/B \\
pGK1025B\_3  & T4SS      & $\sim$16.4 & 15.4 & virB3/4/9/10/11, conj.\ relaxase, phospholipase D \\
\bottomrule
\end{tabular}

\subsection{SSU5 phage homology}
Direct \texttt{blastn} of \texttt{NC\_018843} (\emph{Salmonella} phage SSU5, 104{,}828~bp) against
each candidate plasmid:
\begin{tabular}{lrrl}
\toprule
Plasmid & \# HSPs & Merged qcov (\%) & Note \\
\midrule
pH322\_1 (CP078111)    & 21 & 56.6 & Best HSP 90.2\% id over 5.3~kb \\
pGK1025B\_1 (CP078107) & 22 & 67.0 & Best HSP 91.5\% id over 5.4~kb \\
\bottomrule
\end{tabular}

Both plasmids show large-scale, high-identity SSU5 homology consistent with a phage-plasmid
architecture. PHASTER's ``intact $\sim$96.9~Kbp'' classification scores on phage-gene composition
rather than raw nucleotide coverage, so the two measurements are complementary rather than
identical.

\section{Section-by-section: what worked, what did not}

\subsection{What worked cleanly}
\begin{itemize}
\item \textbf{Data availability.} All seven accessions live at NCBI, deposited as promised
      (BioProject \texttt{PRJNA258403} under GenomeTrakr parent \texttt{PRJNA186875}).
\item \textbf{Sequence identity.} 7/7 exact length match to paper Table 1.
\item \textbf{GC composition.} 7/7 within rounding of paper Table 1 ($\leq 0.18$ pp).
\item \textbf{MLST}: single PubMLST live query returned ST83/CC83 and ST64/CC64 -- exact match.
\item \textbf{PlasmidFinder negativity}: local re-run at CGE defaults gave 0 hits -- confirms
      paper's negative call, with quantitative best-hit-below-threshold data to back it up.
\item \textbf{Annotation-level claims}: T4SS/T6SS/arsenic-operon/phospholipase-D/SSU5 phage all
      present in the current GenBank annotations exactly as the paper describes.
\end{itemize}

\subsection{What partially worked (worth flagging)}
\begin{itemize}
\item \textbf{CDS counts.} Drift $-127..+360$. Root cause: 2026 NCBI PGAP re-annotation vs
      2022 PGAP the paper used. Not a contradiction, but blocks literal reproduction of Table 1's
      CDS numbers without pulling the 2022 archive snapshot. See Open~Question Q1.
\item \textbf{T6SS span.} Our 16--17.5~Kbp outer span vs paper's ``$\sim$13~Kbp truncated core.''
      Definitional difference in which flanking accessory genes count. All key structural genes
      confirmed; span numbers not directly comparable.
\end{itemize}

\subsection{What we did not attempt}
\begin{itemize}
\item \textbf{PHASTER re-run} (paper claims C13, C15, part of C11). PHASTER is webserver-only;
      out of scope for this replication turn. See Q2 for open-source alternative next-step.
\item \textbf{Plasmid phylogeny} (paper claim C14). Multi-plasmid alignment + tree building not
      run; the SSU5-homology signal in isolation is complementary evidence.
\end{itemize}

\section{Verdict}
\textbf{REPLICATED} --- 12/15 claims verified, 0/15 contradicted, 3/15 not attempted (all
webserver-dependent or heavy phylogeny, transparently noted).

\section{Open Questions (Q1--Q5)}\label{sec:openq}
See \texttt{report/open\_questions.json} for the full machine-readable form with per-question
next steps. Titles:
\begin{description}
\item[Q1] CDS-count drift: quantify how much of the paper's 2022 PGAP annotation survived the
  2026 re-annotation, and classify each delta (merge / split / new / removed / pseudogene reclass).
\item[Q2] SSU5 prophage boundaries: use PHASTEST or DBSCAN-SWA CLI to compare the ``intact 96.9
  Kbp'' PHASTER call to our direct BLAST coverage; identify accessory cargo.
\item[Q3] Is \emph{pGK1025B\_3} actually transfer-competent? Biparental mating assay proposal.
\item[Q4] Does the arsenic operon confer phenotypic arsenic tolerance? MIC assay proposal with
  plasmid-cured and arsC-knockout comparators.
\item[Q5] Should CGE PlasmidFinder add a Cronobacter-specific replicon family? Concrete proposal
  based on the 80--91\% best-hit gap we observed.
\end{description}

\section{Failure analysis}\label{sec:failure}
See \texttt{report/failure\_analysis.md} for the full write-up. Summary: CDS-drift is annotation
era (not a paper error), PHASTER was skipped (webserver-only), phylogeny was skipped (heavy and
non-load-bearing), T6SS Kbp is a definitional difference. No hard failures; nothing to work around
beyond retrying the paper-PDF fetch through BMC after NCBI's PMC endpoint returned a pre-fetch
stub.

\section{Data and code availability}
\begin{itemize}
\item Sequences: NCBI GenBank \texttt{CP078106}--\texttt{CP078112} (all live 2026-07-06).
\item Reference phage: \texttt{NC\_018843} (\emph{Salmonella} phage SSU5).
\item CGE PlasmidFinder DB: \url{https://bitbucket.org/genomicepidemiology/plasmidfinder_db}.
\item PubMLST \emph{Cronobacter} seqdef: \url{https://rest.pubmlst.org/db/pubmlst_cronobacter_seqdef}.
\item All replication code and raw outputs live under
      \texttt{\textasciitilde/Dropbox/REPLICATE-PROJECT/BVBRC-128-Negrete-*/work/} and
      \texttt{report/evidence/}.
\end{itemize}

\end{document}
